\section{From Durable State Change to Policy Uptake}
\label{sec:variance}
A different memory string is not yet a different policy. We therefore trace the same reward-label intervention through four distinct gates: durable-state construction, native retrieval exposure, branch-specific policy uptake, and terminal outcome. This ordering is central to the claim: a positive forced intervention can coexist with weak native transport without contradiction.

\subsection{Forced injection establishes latent leverage}
The original action-reach experiment holds a future browser observation fixed while swapping only the paired persistent memory. Seven of 12 aligned rollouts choose different structured next actions, which is an existence witness rather than a population effect. A stronger first-action distribution stress test has mean empirical TV distance 0.344 but permutation $p=0.311$, so systematic action-distribution separation is not established.

The terminal forced memory-swap test fully crosses the four original complete source-memory pairs with four future evidence packets frozen before terminal-policy calls. The initial three-rollout-per-condition run is a committed non-pass (mean $|\Delta|=0.145833$, $p=0.160128$). Before any confirmatory call, the follow-up freezes the identical four sources, four futures, 0.15 practical-effect floor, and $p<0.05$ criterion; it changes only fresh replication depth from three to eight and forbids outcome-driven source or future selection. All 256 confirmatory calls complete, yielding mean $|\Delta|=0.15625$ with permutation $p=0.00074$. Eight of 16 cells are zero, and the two largest cells contain 78.2\% of squared-effect mass. We therefore call this \emph{forced-intervention sensitivity}: the paired states can have conditional leverage when directly injected, but this experiment bypasses native retrieval.

A fresh source-independent no-memory control on these four futures also rules out a simple ``success memory helps / failure memory hurts'' ordering. Omission matches both branches in eight source--future cells, lies closer to success memory in six, and closer to failure memory in two.

\subsection{Native exposure is a separate gate}
ReasoningBank retrieves over task descriptions with all-MiniLM-L6-v2, the current intent as query, top-$1$ selection, and threshold 0.3; reward-conditioned memory content is not part of the retrieval embedding. Exact replay of this released mechanism uses zero provider calls. With only the original four source descriptions, 8 of 188 held-out Shopping tasks clear the threshold (4.26\%), and none of the four forced-terminal futures does. The forced $4\times4$ result is therefore not source-faithful retrieval transport and should not be reported as such.

After adding the 16 preregistered breadth sources, the same top-$1$/0.3 mechanism hits 125 of 172 held-out Shopping tasks (72.7\%), spanning 39 intent-template families. Before new terminal outcomes, we freeze every hit that also has reconstructable released evidence and a deterministic evaluator. This produces 36 native tasks across 13 intent-template families and 11 retriever-selected source memories. Exposure can therefore be broad even when the branch-specific effect that follows exposure is not.

This distinction is consistent with recent trajectory-memory work that treats post-retrieval reuse as a separate problem from selecting a relevant past trajectory~\citep{li2026beyondretrieval}. Our setting differs in treatment---we intervene on reward-conditioned writing---but the evaluation principle is aligned: retrieval success does not determine how strongly the retrieved support changes a policy.

\subsection{Shopping attenuation appears before branch-specific policy uptake}
For each of the 36 frozen native tasks, the retriever itself chooses the source. We recover exact success- and failure-conditioned memory bytes and serialize both through the released ReasoningBank human-memory wrapper. Future evidence, Doubao Seed 2.0 Mini, temperature, evaluator, and four-rollout depth are matched. All 288 calls complete without provider failure. Mean terminal S/F separation is only 0.02083 with permutation $p=0.4289$, far below the unchanged 0.15 practical floor. Thirty-four of 36 tasks have zero branch contrast: 18 joint floors and 16 joint ceilings.

We next test three simple explanations for this attenuation. First, literal omission produces mean memory-presence effect 0.04514 ($p=0.00147$), statistically detectable but below the 0.15 floor. Second, replacing rewritten memory with the exact compact writer-input trajectory yields mean rewrite-versus-raw effect 0.04514 ($p=0.00775$), again below the floor; 31/36 tasks have identical S/F/raw/no-memory point estimates. Third, we hold structured rewriting itself fixed. An outcome-blind procedural writer receives the same task and action-summary bytes, the same DeepSeek-V4-Flash writer family, and the same at-most-three Title/Description/Content schema, but no success/failure, reward, score, or outcome semantics. All 20 neutral memories complete and are textually distinct from both reward branches. On the same 36 native tasks, 144/144 neutral-memory terminal calls complete with zero provider failure. The preregistered reward-conditioned-versus-neutral effect is 0.04514 with $p=0.0048$: detectable, but again below the 0.15 floor. Thirty-two tasks have zero effect and 30/36 have identical success, failure, neutral, raw-trajectory, and no-memory point estimates.

The nonzero neutral-control effect is also highly localized. A zero-call descriptive concentration audit shows that future task 229 alone contributes 61.5\% of absolute effect mass and 87.7\% of squared-effect mass; removing that task lowers the mean effect from 0.04514 to 0.01786. Only two of the 11 retriever-selected source memories have any nonzero future task. This audit does not alter the preregistered permutation result or gate; it clarifies why a small $p$-value should not be read as broad practical transport.

Endpoint resolution is not the main explanation either. A post-hoc, zero-call fractional-reference re-score of 432 archived S/F/no-memory outcomes reduces joint S/F floors from 18 to 10, but mean S/F separation remains 0.01951 overall and 0.02827 on the 16 multi-reference tasks. Because this metric was introduced after observing saturation, it remains diagnostic-only and cannot replace the frozen binary endpoint.

Most importantly, attenuation is already visible before terminal scoring. A preregistered first-action probe freezes the released step-1 browser state on the same 36 native tasks and samples success-memory, failure-memory, and no-memory conditions with four rollouts each. All 432 calls complete with zero provider or unrecoverable parse failures. Mean S/F first-action TV is 0.06944 with permutation $p=0.5801$, below the inherited 0.20 process floor. Only 9/36 states have any S/F TV, and 0/36 change the modal branch action. Removing click indices lowers mean S/F action-family TV to 0.02778. The larger memory-versus-no-memory TV of 0.17014 with six modal shifts is descriptive only because no confirmatory presence gate was preregistered. Zero-call post-hoc diagnostics give only 0.01659 S/F excess in next-goal text and 0.00335 pair-relative branch shift in model-emitted working memory ($p=0.205$, $d_z=0.148$); the latter correlates descriptively with first-action TV ($r=0.464$) but not terminal effect ($r=-0.023$). Together, these diagnostics place the Shopping attenuation before or at \emph{branch-specific policy uptake}; none is a confirmatory mediation test.

\begin{table}[t]
\centering
\scriptsize
\caption{Stage-resolved evidence. Practical-effect floors remain binding when preregistered; a small $p$-value alone is insufficient.}
\label{tab:experiment-ladder}
\begin{tabular}{@{}p{0.27\linewidth}p{0.18\linewidth}p{0.28\linewidth}p{0.17\linewidth}@{}}
\toprule
Scientific gate & Support / calls & Key result & Verdict \\
\midrule
Write identification & 20 complete pairs & 20/20 content/title divergence; pooled Jaccard .673 & support \\
Forced leverage & 16 cells / 256 & mean $|\Delta|=.15625$, $p=.00074$ & pass \\
Native exposure & 172 held-out / 0 & 125 hits (72.7\%); 36 terminal-eligible & support \\
First-action uptake & 36 states / 432 & S/F TV .06944, $p=.5801$; 0/36 modal diff. & non-pass \\
Native terminal & 36 tasks / 288 & S/F $|\Delta|=.02083$, $p=.4289$ & non-pass \\
Outcome-blind structured control & 20 writer + 144 terminal & reward-vs-neutral .04514, $p=.0048$; 32/36 zero & floor fail \\
Reddit cross-domain replication & 4 pairs + 64 terminal & write 4/4; native $|\Delta|=.125$, $p=.2253$ & floor fail \\
\bottomrule
\end{tabular}
\end{table}

\subsection{Generalization and unresolved boundaries}
A second-domain Reddit replication is qualified before any new writer output: eight deterministic native futures across two templates retrieve four source identities. After one uniform output-cap repair that regenerates all eight writer units fresh, all four S/F source pairs diverge (Jaccard 0.652). All 64 terminal calls complete; mean $|\Delta|=0.125$ with $p=0.2253$, below the unchanged 0.15 floor, with six of eight tasks zero and opposite signs in the two nonzero cells. Every leave-one-task-out mean remains below 0.15. Full qualification and repair provenance appears in Appendix~\ref{sec:breadth-transport-details}.

Cross-setting evidence therefore favors a bounded conclusion: write divergence transfers more readily than behavioral magnitude. GLM-5.3 also replicates write divergence but misses its terminal floor (0.140625, $p=0.00012$) with two sign reversals, while the DeepSeek downstream-policy transfer remains a zero-scientific-unit provider-support stop rather than a scientific null. Shopping localizes attenuation before or at branch-specific uptake; Reddit shows that native magnitude itself is domain/task dependent.
